% ============================================
% 数学建模竞赛论文主文件
% 由 mathskills 自动生成与管理
% 编译: xelatex -interaction=nonstopmode 论文.tex (×2)
% ============================================

\documentclass[12pt,a4paper]{ctexart}

% === 页面设置 ===
\usepackage[top=2.5cm, bottom=2.5cm, left=2.8cm, right=2.8cm]{geometry}

% === 基础宏包 ===
\usepackage{etoolbox}
\usepackage{ctex}
\usepackage{url}
\usepackage{graphicx}
\usepackage{amsmath,amssymb}
\usepackage{array}
\usepackage{tabularx}
\usepackage{longtable}
\usepackage{booktabs}
\usepackage{multirow}
\usepackage{enumitem}

% === TikZ 流程图 ===
\usepackage{tikz}
\usetikzlibrary{arrows.meta}
\tikzset{
  box/.style={rectangle, draw, minimum width=3.2cm, minimum height=0.9cm, align=center, font=\small},
  arrow/.style={thick, -{Stealth}}
}

% === 表格列类型 ===
\newcolumntype{C}{>{\centering\arraybackslash}X}
\newcolumntype{R}{>{\raggedleft\arraybackslash}X}
\newcolumntype{L}{>{\raggedright\arraybackslash}X}

% === 标题加粗 ===
\renewcommand{\textbf}[1]{{\song\bfseries #1}}

% === 超链接 ===
\usepackage{hyperref}
\hypersetup{
  colorlinks=true,
  linkcolor=blue!60!black,
  citecolor=blue!40!black,
  urlcolor=blue!60!black
}

% ============================================
\begin{document}

% === 标题 ===
\title{\heiti\zihao{2} 【论文标题——由AI自动填充】}
\author{}
\date{}
\maketitle

% === 摘要 ===
% ============================================
% 0. 摘要
% AI 根据实际赛题内容自动填充
% ============================================

\begin{abstract}

【第一段：开头】本文结合【赛题背景】，研究了【核心研究内容】。首先，【主要工作1】。接着，【主要工作2】。然后，【主要工作3】。最后，【主要工作4】。

【第二段：问题一】\textbf{针对问题一：}本问是\textbf{【问题类型】}问题，需【要求】。首先，对附件数据进行预处理，提取\textbf{【关键特征】}。接着，采用\textbf{【主要方法/模型】}进行\textbf{【建模/求解过程】}。结果显示，\textbf{【具体数值结果】}，模型有效地\textbf{【评价】}。

【第三段：问题二】\textbf{针对问题二：}本问是\textbf{【问题类型】}问题，需【要求】。通过分析\textbf{【机理/数据特征】}，建立了\textbf{【模型名称】}。利用\textbf{【求解方法】}求解得到\textbf{【具体数值结果】}，结果验证了\textbf{【结论】}。

【第四段：问题三】\textbf{针对问题三：}本问是\textbf{【问题类型】}问题，需【要求】。基于\textbf{【前文结果/新模型】}，通过\textbf{【方法】}进行求解，得到\textbf{【具体数值结果】}，为\textbf{【决策/分析】}提供了依据。

【第五段：问题四（如有）】\textbf{针对问题四：}本问是\textbf{【问题类型】}问题，需【要求】。从\textbf{【角度】}出发，给出\textbf{【建议/结论】}。

【第六段：结尾】针对模型的可靠性，本文进行了\textbf{【灵敏度分析方法】}和\textbf{【交叉验证方法】}，验证了模型的\textbf{【稳健性和准确性】}。综上，本文为\textbf{【领域】}问题提供了\textbf{【高效/可靠】}的解决方案。

\keywords{【关键词1】；【关键词2】；【关键词3】；【关键词4】；【关键词5】；【关键词6】}
\label{abstract:end}

\end{abstract}


% === 目录 ===
\tableofcontents
\newpage

% === 正文 ===
% ============================================
% 1. 问题重述
% ============================================

\section{问题重述}

\subsection{问题背景}

【宏观背景，约300字，需引用2-3篇参考文献】

随着\textbf{【领域背景】}的快速发展，\textbf{【具体问题】}日益受到关注。近年来，学者们围绕该问题开展了大量研究。\cite{ref1} 提出了\textbf{【已有方法1】}，该方法在\textbf{【场景】}中表现良好，但存在\textbf{【局限性1】}。\cite{ref2} 通过\textbf{【已有方法2】}对\textbf{【问题】}进行改进，然而\textbf{【局限性2】}仍未得到有效解决。\cite{ref3} 的研究表明，\textbf{【关键发现】}，这为解决该问题提供了新的思路。

当前，\textbf{【研究现状总结】}。因此，亟需一种\textbf{【本文要解决的问题】}的方法，以\textbf{【研究目标】}。

\subsection{问题提出}

本文所研究的问题来源于\textbf{【赛题来源】}。附件数据包含\textbf{【数据描述：文件类型、数据量、关键字段】}，涵盖了\textbf{【数据覆盖的时间范围/对象/维度】}等信息。

基于附件数据，本文需要解决以下问题：

\textbf{问题一：}【问题一的完整描述，包括输入、输出和约束条件】

\textbf{问题二：}【问题二的完整描述，包括输入、输出和约束条件】

\textbf{问题三：}【问题三的完整描述，包括输入、输出和约束条件】

\textbf{问题四：}【问题四的完整描述，包括输入、输出和约束条件】

% ============================================
% 2. 问题分析
% 小标题格式：2.N 问题N——基于XX的XX分析
% ============================================

\section{问题分析}

\subsection{问题一——基于【主要方法】的【问题类型】分析}

针对问题一，该问题的本质是一个\textbf{【问题类型：预测/分类/评价/优化】}问题。考虑到附件数据具有\textbf{【数据特征：高维度/时序性/非线性/多变量】}的特点，拟采用\textbf{【主要模型/方法】}进行求解。

分析流程如下：首先，对附件数据进行\textbf{【预处理方式】}，提取\textbf{【关键特征】}；接着，建立\textbf{【模型名称】}模型，通过\textbf{【参数确定方法】}确定模型参数；然后，利用\textbf{【求解算法】}进行求解；最后，对模型结果进行\textbf{【验证方式】}和\textbf{【灵敏度分析】}。

具体分析流程如图\ref{fig:analysis1_flow}所示。

\begin{figure}[ht]
  \begin{center}
    \begin{tikzpicture}[x=1cm, y=0.7cm]
      \node[box] (n11) at (0,0) {数据预处理};
      \node[box] (n12) at (5,0) {特征提取};
      \node[box] (n13) at (10,0) {模型构建};
      \node[box] (n22) at (5,-3) {模型求解};
      \node[box] (n21) at (0,-3) {结果验证};
      \node[box] (n23) at (10,-3) {参数优化};
      \node[box] (n31) at (0,-6) {结果输出};
      \node[box] (n32) at (5,-6) {灵敏度分析};
      \node[box] (n33) at (10,-6) {结论总结};
      \draw[arrow] (n11)--(n12); \draw[arrow] (n12)--(n13);
      \draw[arrow] (n23)--(n22); \draw[arrow] (n22)--(n21);
      \draw[arrow] (n13) -- ++(0,-1.0) -| (n23);
      \draw[arrow] (n21)--(n31); \draw[arrow] (n31)--(n32); \draw[arrow] (n32)--(n33);
    \end{tikzpicture}
    \caption{问题一分析流程图}
    \label{fig:analysis1_flow}
  \end{center}
\end{figure}

\subsection{问题二——基于【主要方法】的【问题类型】分析}

针对问题二，该问题的本质是一个\textbf{【问题类型】}问题。考虑到\textbf{【问题特点】}，拟采用\textbf{【主要模型/方法】}进行求解。

分析流程如下：首先，\textbf{【步骤1】}；接着，\textbf{【步骤2】}；然后，\textbf{【步骤3】}；最后，\textbf{【步骤4】}。

具体分析流程如图\ref{fig:analysis2_flow}所示。

\begin{figure}[ht]
  \begin{center}
    \begin{tikzpicture}[x=1cm, y=0.7cm]
      \node[box] (n11) at (0,0) {...};
      \node[box] (n12) at (5,0) {...};
      \node[box] (n13) at (10,0) {...};
      \node[box] (n22) at (5,-3) {...};
      \node[box] (n21) at (0,-3) {...};
      \node[box] (n23) at (10,-3) {...};
      \node[box] (n31) at (0,-6) {...};
      \node[box] (n32) at (5,-6) {...};
      \node[box] (n33) at (10,-6) {...};
      \draw[arrow] (n11)--(n12); \draw[arrow] (n12)--(n13);
      \draw[arrow] (n23)--(n22); \draw[arrow] (n22)--(n21);
      \draw[arrow] (n13) -- ++(0,-1.0) -| (n23);
      \draw[arrow] (n21)--(n31); \draw[arrow] (n31)--(n32); \draw[arrow] (n32)--(n33);
    \end{tikzpicture}
    \caption{问题二分析流程图}
    \label{fig:analysis2_flow}
  \end{center}
\end{figure}

\subsection{问题三——基于【主要方法】的【问题类型】分析}

针对问题三，该问题的本质是一个\textbf{【问题类型】}问题。考虑到\textbf{【与前序问题的递进关系】}，拟采用\textbf{【主要模型/方法】}进行求解。

分析流程如下：首先，\textbf{【步骤1】}；接着，\textbf{【步骤2】}；然后，\textbf{【步骤3】}；最后，\textbf{【步骤4】}。

\subsection{问题四——基于【主要方法】的【问题类型】分析}

针对问题四，该问题的本质是一个\textbf{【问题类型】}问题。基于\textbf{【前三个问题的分析结果】}，拟从\textbf{【角度1】}和\textbf{【角度2】}出发，给出\textbf{【建议/改进方案/推广方向】}。

% ============================================
% 3. 模型假设
% ============================================

\section{模型假设}

为简化问题，本文做出以下假设：

\begin{itemize}
  \item \textbf{假设1：}【假设内容】。合理性说明：在\textbf{【实际场景】}中，\textbf{【为什么合理】}，该假设对模型精度的影响在可接受范围内。

  \item \textbf{假设2：}【假设内容】。合理性说明：根据\textbf{【文献/数据观察】}，\textbf{【为什么合理】}。

  \item \textbf{假设3：}【假设内容】。合理性说明：考虑到\textbf{【实际约束】}，\textbf{【为什么合理】}。

  \item \textbf{假设4：}【假设内容】。合理性说明：在\textbf{【研究范围内】}，\textbf{【为什么合理】}。
\end{itemize}

% ============================================
% 4. 符号说明 —— 三线表格式
% ============================================

\section{符号说明}

本文中所使用的主要符号及其说明如表\ref{tab:symbols}所示。

\begin{longtable}{>{\centering\arraybackslash}p{0.22\textwidth}>{\centering\arraybackslash}p{0.60\textwidth}>{\centering\arraybackslash}p{0.10\textwidth}}
  \caption{符号说明表}
  \label{tab:symbols} \\
  \toprule
  \textbf{符号} & \textbf{说明} & \textbf{单位} \\
  \midrule
  \endfirsthead
  \caption{符号说明表（续）} \\
  \toprule
  \textbf{符号} & \textbf{说明} & \textbf{单位} \\
  \midrule
  \endhead
  \bottomrule
  \endfoot
  $x_i$ & 第$i$个样本的特征值 & — \\
  $N$ & 样本总数 & 个 \\
  $y_i$ & 第$i$个样本的目标值 & — \\
  $\hat{y}_i$ & 第$i$个样本的预测值 & — \\
  $\bar{x}$ & 样本均值 & — \\
  $\sigma$ & 样本标准差 & — \\
  $R^2$ & 决定系数 & — \\
  $w_j$ & 第$j$个评价指标的权重 & — \\
  % 【AI 根据实际问题自动补充符号】
\end{longtable}

% ============================================
% 5. 问题一的模型构建与求解
% ============================================

\section{问题一的模型构建与求解}

\subsection{求解思路}

针对问题一，本问是一个\textbf{【问题类型：预测/分类/评价/优化】}问题。本文首先对附件数据进行\textbf{【预处理】}，提取\textbf{【关键特征】}；接着，通过算法对比选择\textbf{【最优算法】}；然后，建立\textbf{【模型名称】}模型并采用\textbf{【求解方法】}进行求解；最后，对结果进行\textbf{【验证与分析】}。

\subsection{模型建立}

\textbf{Step 1：数据预处理}

（1）预处理原因：原始数据来源于\textbf{【赛题附件】}，格式为\textbf{【xlsx/csv】}。原始数据存在\textbf{【缺失值/量纲不统一/异常值/多级表头】}等问题，必须进行预处理以确保建模质量。

（2）数据理解与提取：从附件数据中提取\textbf{【关键字段列表】}，共计\textbf{【N】}个特征、\textbf{【M】}条样本。

（3）预处理步骤：
\begin{itemize}
  \item 数据清洗——对缺失值采用\textbf{【均值/中位数/线性插值】}填充；对异常值采用\textbf{【IQR法 / Z-score法（阈值=3）】}检测并替换为\textbf{【边界值】}。
  \item 标准化——为消除量纲影响，采用\textbf{【Min-Max归一化 / Z-score标准化】}：
  \begin{equation}
  x' = \frac{x - x_{\min}}{x_{\max} - x_{\min}}
  \end{equation}
  其中，$x$为原始值，$x_{\min}$为最小值，$x_{\max}$为最大值，$x'$为归一化后的值。
  \item \textbf{【如有特征衍生，在此展开】}
\end{itemize}

（4）预处理结果：处理前后关键统计量对比如表\ref{tab:preprocess}所示。

\begin{longtable}{>{\centering\arraybackslash}p{0.12\textwidth}>{\centering\arraybackslash}p{0.17\textwidth}>{\centering\arraybackslash}p{0.17\textwidth}>{\centering\arraybackslash}p{0.17\textwidth}>{\centering\arraybackslash}p{0.17\textwidth}}
  \caption{数据预处理前后统计量对比}
  \label{tab:preprocess} \\
  \toprule
  指标 & 处理前均值 & 处理前标准差 & 处理后均值 & 处理后标准差 \\
  \midrule \endfirsthead
  \caption{数据预处理前后统计量对比（续）} \\ \toprule
  指标 & 处理前均值 & 处理前标准差 & 处理后均值 & 处理后标准差 \\ \midrule \endhead
  \bottomrule \endfoot
  ... & ... & ... & ... & ... \\
\end{longtable}

（5）预处理完成，数据已规整，为后续建模提供了清洁输入。

\textbf{Step 2：算法选择}

该问题本质上是一个\textbf{【问题类型】}问题，需要在\textbf{【条件】}下寻找\textbf{【目标】}。常用求解算法包括\textbf{【算法A】}、\textbf{【算法B】}、\textbf{【算法C】}等。本节选取\textbf{【3-4】}种经典算法，从\textbf{【维度1】}、\textbf{【维度2】}、\textbf{【维度3】}、\textbf{【维度4】}等方面进行评估，结果如表\ref{tab:algo_compare}所示。

\begin{longtable}{>{\centering\arraybackslash}p{0.12\textwidth}>{\centering\arraybackslash}p{0.18\textwidth}>{\centering\arraybackslash}p{0.18\textwidth}>{\centering\arraybackslash}p{0.18\textwidth}>{\centering\arraybackslash}p{0.18\textwidth}}
  \caption{候选算法能力对比}
  \label{tab:algo_compare} \\
  \toprule
  算法名称 & 全局搜索能力 & 收敛速度 & 约束处理能力 & 求解稳定性 \\
  \midrule \endfirsthead
  \caption{候选算法能力对比（续）} \\ \toprule
  算法名称 & 全局搜索能力 & 收敛速度 & 约束处理能力 & 求解稳定性 \\ \midrule \endhead
  \bottomrule \endfoot
  遗传算法（GA） & 优 & 优 & 良 & 良 \\
  蚁群算法（ACO） & 良 & 中 & 优 & 中 \\
  模拟退火（SA） & 中 & 良 & 良 & 中 \\
\end{longtable}

通过上述对比可以看出，\textbf{【选定算法（英文缩写）】}在\textbf{【优势维度】}方面表现最优，因此本文选择该算法作为核心求解方法。

\textbf{Step 3：模型构建}

本文采用\textbf{【模型全称】（【英文缩写】）}建立数学模型。该模型的基本形式为：

\begin{equation}
y = f(\mathbf{x}; \boldsymbol{\theta})
\end{equation}

其中，$y$为\textbf{【输出/预测目标】}，$\mathbf{x}$为\textbf{【输入特征向量】}，$\boldsymbol{\theta}$为\textbf{【模型参数】}。

% === AI 在此继续展开公式推导，每步≥6-7公式 ===

\textbf{Step 4：模型求解}

采用\textbf{【求解算法】}对上述模型进行求解。参数设定为：\textbf{【参数1】}=\textbf{【值】}（来源：\textbf{【网格搜索/文献/数据统计】}），\textbf{【参数2】}=\textbf{【值】}。

\subsection{结果分析}

求解结果如表\ref{tab:result1}所示。

\begin{longtable}{>{\centering\arraybackslash}p{0.20\textwidth}>{\centering\arraybackslash}p{0.20\textwidth}>{\centering\arraybackslash}p{0.20\textwidth}>{\centering\arraybackslash}p{0.20\textwidth}}
  \caption{问题一求解结果}
  \label{tab:result1} \\
  \toprule
  指标1 & 指标2 & 指标3 & 指标4 \\
  \midrule \endfirsthead
  \caption{问题一求解结果（续）} \\ \toprule
  指标1 & 指标2 & 指标3 & 指标4 \\ \midrule \endhead
  \bottomrule \endfoot
  ... & ... & ... & ... \\
\end{longtable}

结果可视化如图\ref{fig:result1}所示。

\begin{figure}[ht]
  \centering
  \includegraphics[width=0.8\textwidth]{../../求解/问题一/图片/【图片名】.png}
  \caption{问题一结果可视化}
  \label{fig:result1}
\end{figure}

如图\ref{fig:result1}所示，\textbf{【对图表的具体分析，≥3行】}。从图中可以观察到，\textbf{【趋势/规律/异常点】}。这一结果表明\textbf{【与题目要求的关联】}。

综上所述，问题一的求解结果表明，\textbf{【核心结论】}。模型的\textbf{【精度指标】}达到了\textbf{【具体数值】}，验证了所建模型的有效性。

% ============================================
% 6. 问题二的模型构建与求解
% ============================================

\section{问题二的模型构建与求解}

\subsection{求解思路}

针对问题二，本问是一个\textbf{【问题类型】}问题。基于问题一的\textbf{【结果/预处理数据】}，本文首先\textbf{【步骤1】}；接着，建立\textbf{【模型名称】}模型；然后，通过\textbf{【求解方法】}求解；最后，对结果进行\textbf{【验证与分析】}。

\subsection{模型建立}

\textbf{Step 1：【第一步标题】}

……

\begin{equation}
% 公式内容
\end{equation}

其中，$X$表示\textbf{【变量含义】}。

\textbf{Step 2：【第二步标题】}

……

\subsection{结果分析}

……（图表 + 分析）……

% ============================================
% 7. 问题三的模型构建与求解
% ============================================

\section{问题三的模型构建与求解}

\subsection{求解思路}

针对问题三，本问是一个\textbf{【问题类型】}问题。基于\textbf{【前序问题结果】}，本文首先\textbf{【步骤1】}；接着，\textbf{【步骤2】}；然后，\textbf{【步骤3】}；最后，\textbf{【步骤4】}。

\subsection{模型建立}

% AI 自动填充
……

\subsection{结果分析}

% AI 自动填充
……

% ============================================
% 8. 问题四的模型构建与求解
% ============================================

\section{问题四的模型构建与求解}

\subsection{求解思路}

针对问题四，本问是一个\textbf{【问题类型】}问题。基于\textbf{【前三个问题的分析结果】}，本文从\textbf{【角度1】}和\textbf{【角度2】}出发，给出\textbf{【建议/改进方案】}。

\subsection{分析与建议}

% AI 自动填充（建议总结型问题以自然段落为主）
……

\subsection{结果分析}

% AI 自动填充
……

% ============================================
% 9. 模型灵敏度分析
% ============================================

\section{模型灵敏度分析}

为检验模型的\textbf{稳健性和可靠性}，本文对关键参数进行了灵敏度分析。通过改变\textbf{【关键参数列表】}等参数的取值，观察其对模型结果的\textbf{【影响程度】}。

\subsection{问题一的灵敏度分析}

针对问题一中的\textbf{【关键参数1】}和\textbf{【关键参数2】}，在\textbf{【原始值 $\pm 20\%$】}的范围内进行灵敏度测试，结果如图\ref{fig:sensitivity1}所示。

\begin{figure}[ht]
  \centering
  \includegraphics[width=0.8\textwidth]{../../求解/问题一/图片/【灵敏度分析图】.png}
  \caption{问题一参数灵敏度分析}
  \label{fig:sensitivity1}
\end{figure}

如图\ref{fig:sensitivity1}所示，当\textbf{【参数】}在\textbf{【范围】}内变化时，模型结果的\textbf{【变化趋势】}。具体而言，\textbf{【具体分析，≥3行】}。这表明模型对\textbf{【参数名】}的敏感程度\textbf{【高/低】}，模型的\textbf{【稳健性评价】}。

\subsection{问题二的灵敏度分析}

……

% 每个问题都必须配灵敏度分析图和 ≥3 行分析文字

% ============================================
% 10. 模型评价与推广
% ============================================

\section{模型评价与推广}

\subsection{模型优点}

\begin{itemize}
  \item \textbf{优点1：模型精度高。}本模型在\textbf{【数据集/场景】}上取得了\textbf{【具体指标数值】}的优异表现，相较于\textbf{【对比方法】}提升了\textbf{【百分比】}，验证了模型的\textbf{【准确性/有效性】}。

  \item \textbf{优点2：模型可解释性强。}通过\textbf{【SHAP分析/特征重要性排序/公式显式表达】}，模型的\textbf{【决策过程/预测依据】}具有清晰的可解释性，便于\textbf{【实际应用/决策支持】}。

  \item \textbf{优点3：模型鲁棒性好。}灵敏度分析结果表明，模型在\textbf{【参数变动范围】}内表现稳定，核心结论不受参数扰动的影响，具有良好的\textbf{【泛化能力/稳健性】}。

  \item \textbf{优点4：建模方法系统完整。}本文从\textbf{【数据处理】}到\textbf{【模型建立】}再到\textbf{【结果验证】}，形成了完整的分析流程，为\textbf{【同类问题】}提供了可参考的范式。
\end{itemize}

\subsection{模型缺点}

\begin{itemize}
  \item \textbf{缺点1：数据依赖性。}模型的性能依赖于\textbf{【数据质量/数据量】}，在\textbf{【数据缺失严重/数据分布变化】}的情况下，模型精度可能\textbf{【下降】}。未来可通过\textbf{【数据增强/迁移学习】}等方法改善。

  \item \textbf{缺点2：计算复杂度。}模型在\textbf{【训练/求解】}过程中需要\textbf{【较多计算资源/较长计算时间】}，对于\textbf{【大规模/实时】}应用场景可能存在\textbf{【效率瓶颈】}。可通过\textbf{【算法优化/并行计算】}加速。
\end{itemize}

\subsection{模型推广}

本文所建立的模型不仅适用于\textbf{【当前赛题场景】}，还可推广至\textbf{【相关领域/更广泛场景】}。例如，在\textbf{【领域A】}中，可将\textbf{【模型核心方法】}应用于\textbf{【具体问题】}；在\textbf{【领域B】}中，通过\textbf{【调整参数/扩展维度】}即可解决\textbf{【类似问题】}。此外，模型的\textbf{【核心思想/技术框架】}对于\textbf{【更广泛的问题类型】}也具有一定的借鉴意义。


% === 参考文献（单独一页） ===
\newpage
% ============================================
% 参考文献（单独一页）
% GB/T 7714-2015 格式，8-15 条
% AI 根据论文实际引用内容自动生成
% ============================================

\begin{thebibliography}{99}

\bibitem{ref1} 作者1, 作者2. 论文标题[J]. 期刊名, 年份, 卷(期): 起止页码.

\bibitem{ref2} 作者. 论文标题[C]// 会议名称. 会议地点: 出版社, 年份: 起止页码.

\bibitem{ref3} 作者. 书名[M]. 出版地: 出版社, 年份.

\bibitem{ref4} Author A, Author B. Title of the paper[J]. Journal Name, Year, Volume(Issue): Pages.

\bibitem{ref5} Author C. Title of the book[M]. Place: Publisher, Year.

\bibitem{ref6} 作者3, 作者4. 论文标题[J]. 期刊名, 年份, 卷(期): 页码.

\bibitem{ref7} Author D, Author E. Conference paper title[C]// Conference Name. Location, Year: Pages.

\bibitem{ref8} 作者5. 论文标题[D]. 学校所在地: 学校名称, 年份.

\end{thebibliography}


\end{document}
